\part{Foundational Principles of Data-Driven Curriculum Stratification for Physical Systems}

The training of Physics-Informed Neural Networks (PINNs) on complex, chaotic dynamical systems represents a significant challenge at the intersection of scientific computing and machine learning. A key strategy for improving training stability and final model accuracy is Curriculum Learning (CL), where the model is exposed to training examples in a structured sequence, typically from "easy" to "hard".[1, 2] This report presents a comprehensive, revised workflow for a data-driven curriculum stratification framework tailored to a PINN project for the three-body problem. The proposed methodology moves beyond static, heuristic-based criteria, which have proven brittle, towards an adaptive system that derives its structure from the empirical properties of the simulation dataset itself. This approach is grounded in best practices from machine learning and is designed to be robust, interpretable, and directly implementable.

\section{The Rationale for Adaptive Curricula in Physics Simulations}

The core premise of curriculum learning is to create a smooth "difficulty gradient" that guides the optimization process of a neural network towards a robust and generalizable solution.[3, 4] In the context of physics simulations, difficulty is often tied to the complexity of the underlying dynamics, such as the transition from regular, periodic motion to chaotic evolution. The initial approach for the CTL-PINN project, as detailed in the project documentation, relied on defining curriculum stages using fixed, theoretical thresholds on key physical invariants.[5] For instance, Stage 0 was defined by strict bounds on total energy ($E$), angular momentum ($L/L_{\text{max}}$), virial ratio ($Q$), and minimum particle separation ($s_{\text{min}}$) to isolate highly stable, hierarchical systems.[5]

While physically motivated, this static, heuristic-based approach possesses an inherent fragility: it presupposes that the training data will conform to idealized statistical distributions of these invariants. This assumption is frequently violated in practice, particularly when dealing with complex data generation processes. The project's use of Sobol-based quasi-random sampling for initial conditions (ICs) introduces a scaling factor, $s$, which systematically distorts the virial ratio, since $Q \propto 2T/|V|$ scales approximately as $s^3$.[5] Consequently, a large fraction of the quasi-random simulations, though potentially stable, fall outside the predefined $Q$ range (e.g., $[0.9, 1.1]$ for Stage 3), leading to underpopulated or even empty stages.[5] This mismatch between the curriculum's assumptions and the data's reality undermines the entire CL strategy, creating abrupt jumps in difficulty rather than a smooth progression.

This experience aligns with a broader consensus in machine learning research: "empirical adaptation to dataset statistics outperforms hand-crafted heuristics, especially in noisy or imbalanced data".[5] For PINNs applied to chaotic systems, this principle is especially pertinent. The "noise" is not merely measurement error but the intrinsic, high-frequency dynamics of chaos, and the "imbalance" is the often non-uniform distribution of stability across the phase space. Research on PINN training pathologies confirms that managing complexity is paramount for achieving convergence, with curriculum-based strategies being a primary tool for this purpose.[6, 7, 8] Papers on CL for PINNs advocate for approaches that progressively increase the difficulty, whether by sequencing the data samples or by gradually introducing more complex physical constraints.[4, 7]

This necessitates a paradigm shift. Instead of viewing the curriculum as a static, a priori set of rules, it should be conceptualized as a simplified, data-driven \textit{model of the dataset's difficulty landscape}. A good curriculum, like any good model, must accurately represent the underlying data distribution, capturing its unique modes, skews, and correlations. A curriculum built on faulty assumptions is a mismatched model, leading to poor "predictions" about which data samples are genuinely easy or hard for the PINN to learn. Therefore, the design of the curriculum must be treated with the same analytical rigor as the design of the neural network itself. It requires a dedicated phase of exploratory data analysis (EDA), feature engineering (selection and transformation of physical invariants), model selection (e.g., choosing between clustering and quantile-based methods), and, ultimately, empirical validation. By adopting this data-driven philosophy, the curriculum becomes a robust and adaptive framework that is resilient to the specific artifacts of the data generation process, ensuring a smooth and effective learning path for the PINN.

\section{A Lexicon of Difficulty Metrics for the Three-Body Problem}

To construct a data-driven curriculum, a set of quantitative features that correlate with learning difficulty is required. For the three-body problem, the conserved quantities and geometric properties of the system serve as excellent, physically interpretable difficulty metrics.[5] Each invariant provides a proxy for a specific aspect of dynamical complexity, forming a multi-dimensional feature space from which the curriculum can be constructed. The following metrics, computed from the initial state of each simulation, are foundational to the proposed framework.

\begin{description}
    \item[Total Mechanical Energy ($E$):] Defined as the sum of the system's kinetic ($T$) and potential ($V$) energies.
    $$E = T + V = \frac{1}{2}\sum_{i=1}^{3}m_{i}\lVert v_{i}'\rVert^{2} - \sum_{1\le i<j\le3}\frac{m_{i}m_{j}}{\lVert r_{i}'-r_{j}'\rVert+\epsilon}$$
    In a system with gravitational constant $G=1$, a negative total energy ($E < 0$) indicates a gravitationally bound system. The more negative the value of $E$, the more tightly bound the system is, which often corresponds to more stable, regular orbits. Systems with energy close to zero ($E \approx 0$) are marginally bound or unbound and are susceptible to chaotic scattering, ejections, and fly-by encounters, which are significantly harder for a PINN to model accurately.[5]

    \item[Virial Ratio ($Q$):] The ratio of twice the kinetic energy to the magnitude of the potential energy.
    $$Q = \frac{2T}{|V|}$$
    For a stable, self-gravitating system in statistical equilibrium, the virial theorem states that $Q \approx 1$. A significant deviation from this value implies a state of disequilibrium. A system with $Q > 1$ has excess kinetic energy and will tend to expand, while a system with $Q < 1$ is dominated by potential energy and will tend to contract. Both scenarios often precede chaotic evolution, making $Q$ a key indicator of initial stability.[5]

    \item[Normalized Angular Momentum ($L_{\text{norm}}$):] The magnitude of the total angular momentum vector, $L$, normalized by a characteristic maximum angular momentum, $L_{\text{max}}$.
    $$L_{\text{norm}} = \frac{\lVert L\rVert}{L_{\text{max}}} = \frac{\left\lVert\sum_{i=1}^{3}m_{i}(r_{i}'\times v_{i}')\right\rVert}{M_{\text{tot}}r_{\text{char}}v_{\text{char}}}$$
    where $r_{\text{char}}$ is a characteristic length scale and $v_{\text{char}}$ is a characteristic velocity. A value of $L_{\text{norm}}$ close to 1 corresponds to near-circular orbits, which are highly regular and stable. Lower values indicate more eccentric (elongated) orbits, which lead to periodic close encounters between bodies, a primary driver of chaotic instabilities.[5]

    \item[Minimum Initial Separation ($s_{\text{min}}$):] The smallest Euclidean distance between any pair of bodies in the initial configuration.
    $$s_{\text{min}} = \min_{1\le i<j\le3}\lVert r_{i}'-r_{j}'\rVert$$
    This is a direct geometric indicator of the immediate risk of a close encounter. Very small values of $s_{\text{min}}$ signal impending strong gravitational interactions, which can lead to violent, unpredictable scattering events. These scenarios are among the most numerically challenging for both traditional integrators and PINNs.[5]

    \item[Lyapunov Time ($\tau_L$):] The inverse of the largest Lyapunov exponent ($\lambda_{\text{max}}$), estimated from the simulated trajectory.
    $$\tau_{L} = \frac{1}{\lambda_{\text{max}}}$$
    $\tau_{L}$ quantifies the timescale over which predictability is lost. It is the most direct measure of chaos in a dynamical system, representing the e-folding time for the exponential divergence of initially nearby trajectories. A short $\tau_{L}$ indicates strong chaos and high sensitivity to initial conditions, making long-term prediction fundamentally difficult. A long or infinite $\tau_{L}$ signifies regular, predictable motion.[5]

    \item[Shape Invariants ($I_1, I_2, I_3$):] Dimensionless ratios of the three mutual distances between the bodies ($d_{\text{short}} \le d_{\text{mid}} \le d_{\text{long}}$).
    $$I_{1}=\frac{d_{\text{short}}}{d_{\text{long}}}, \quad I_{2}=\frac{d_{\text{mid}}}{d_{\text{long}}}, \quad I_{3}=\frac{d_{\text{short}}+d_{\text{mid}}-d_{\text{long}}}{d_{\text{short}}+d_{\text{mid}}+d_{\text{long}}}$$
    These invariants describe the geometry of the initial configuration (e.g., collinear, isosceles, equilateral) independent of scale or orientation. Certain initial geometries are known to be more prone to chaotic evolution than others, making these invariants potentially useful secondary features for characterizing difficulty.[5]
\end{description}

These invariants form the feature vector for each simulation. By analyzing the statistical distributions and correlations of these features across the entire dataset, it is possible to construct a curriculum that is quantitatively grounded in the physics of the problem and empirically adapted to the data at hand.

\begin{table}[H]
\centering
\caption{Physical Invariants and Their Interpretation as Difficulty Metrics}
\label{tab:invariants}
\begin{tabular}{@{}p{2cm}p{4cm}p{7cm}@{}}
\toprule
\textbf{Invariant} & \textbf{Full Name} & \textbf{Role as a Difficulty Metric (Easy vs. Hard)} \\ \midrule
$E$ & Total Mechanical Energy & \textbf{Easy:} Large negative values (deeply bound). \textbf{Hard:} Values near zero. \\
$Q_{\text{norm}}$ & Normalized Virial Ratio & \textbf{Easy:} Values near 1 (virial equilibrium). \textbf{Hard:} Values far from 1. \\
$L_{\text{norm}}$ & Normalized Angular Momentum & \textbf{Easy:} Values near 1 (near-circular orbits). \textbf{Hard:} Low values. \\
$s_{\text{min}}$ & Minimum Initial Separation & \textbf{Easy:} Large values (well-separated). \textbf{Hard:} Small values. \\
$\tau_{L}$ & Lyapunov Time & \textbf{Easy:} Large or infinite values (regular). \textbf{Hard:} Small values (strong chaos). \\ \bottomrule
\end{tabular}
\end{table}

\section{Core Techniques for Data-Driven Stratification}

With a well-defined feature space of physical invariants, the next step is to select appropriate machine learning techniques to partition this space into curriculum stages. A hybrid approach combining dimensionality reduction, unsupervised clustering, and quantile-based binning provides a powerful and flexible toolkit for this task.

\subsection{Dimensionality Reduction with PCA for a Composite Stability Score}

In the three-body problem, difficulty is a multi-faceted concept. A system can be challenging due to low angular momentum, near-zero energy, or a short Lyapunov time. These metrics are often strongly correlated; for example, a very tightly bound system (highly negative $E$) may evolve into chaos more quickly (short $\tau_L$).[5] Simply setting independent thresholds on each metric can lead to conflicting or redundant criteria.

Principal Component Analysis (PCA) offers a principled method to address this. PCA is an orthogonal linear transformation that projects the data onto a new coordinate system, where the axes (principal components, PCs) are ordered by the amount of variance they explain in the original data.[9, 10] The first principal component (PC1) represents the direction of maximum variance in the feature space. In the context of our physical invariants, this component can often be interpreted as a composite "stability" or "difficulty" score.[5]

For instance, if PC1 is found to be a linear combination with positive weights on $L_{\text{norm}}$ and $\tau_L$ and a negative weight on $|E|$, then a high score on PC1 would correspond to a system that is simultaneously regular, predictable, and not excessively bound---a clear signature of stability. By ranking simulations based on their PC1 score, one can create a continuous, one-dimensional difficulty axis that optimally combines the information from multiple correlated invariants. This data-driven approach is superior to manually constructing a composite score (e.g., $-E + L_{\text{norm}}$ [5]), as the weights are derived directly from the dataset's covariance structure, ensuring they are optimally tailored to the specific data distribution.[11, 12]

\subsection{Unsupervised Clustering for Discovering Dynamical Regimes}

While PCA and quantile binning are excellent for ordering data along a continuous difficulty axis, they may not capture the full picture. Difficulty is not always just a matter of degree; it can also be a matter of kind. Unsupervised clustering algorithms, such as K-Means or Gaussian Mixture Models (GMM), partition the dataset into distinct groups based on similarity in the multi-dimensional feature space. This allows for the automatic discovery of different \textit{dynamical regimes} or \textit{phenomenological classes} of simulations.[5]

Consider two simulations that have the same short Lyapunov time, $\tau_L$. From a purely quantile-based perspective, they are equally difficult. However, their underlying dynamics could be completely different. One might be a "prompt ejection" event, where one body is quickly flung out of the system. The other might be a long-lived, complex "resonant encounter," where the bodies undergo multiple close approaches before settling into a new configuration. For a PINN, these two scenarios present qualitatively different challenges. The prompt ejection involves large velocities and an unbound trajectory, while the resonant encounter involves intricate, repeated interactions in a compact region of space.

A simple easy-to-hard curriculum might mix these two types of chaos, potentially confusing the learning process. Clustering in the full feature space (e.g., using $E, L_{\text{norm}}, s_{\text{min}}, \tau_L$) can distinguish between them. It might naturally group all prompt ejections into one cluster (characterized, perhaps, by $E \approx 0$ and high final velocities) and all resonant encounters into another (characterized by small $s_{\text{min}}$ and oscillating inter-particle distances). This reveals that the difficulty landscape is not a simple line but a structured space with distinct regions.

This allows for the design of a more sophisticated curriculum that introduces not just progressively harder simulations, but also qualitatively different types of physical phenomena stage by stage. For example, after learning stable orbits, the PINN could be trained specifically on the "prompt ejection" cluster, then on the "resonant encounter" cluster. This aligns with research suggesting that clustering can reveal "latent difficulty gradients" that hand-picked thresholds might overlook.[5, 13] K-Means is an efficient baseline for this task, while GMMs offer more flexibility by modeling non-spherical clusters and providing probabilistic assignments, which can be invaluable for identifying outliers or creating "soft" curriculum boundaries.[5, 14]

\subsection{Quantile-Based Binning for Robust and Balanced Pacing}

Quantile-based binning is a non-parametric method that divides a ranked dataset into a predetermined number of continuous intervals ("bins"), each holding the same number of data points.[15] For example, quartiles divide the data into four bins, each containing 25\% of the samples.

This technique is exceptionally robust and serves as a cornerstone of the adaptive curriculum. Its primary advantage is that it makes no assumptions about the underlying distribution of the difficulty metric. This is critical for the CTL-PINN project, where metrics like the raw virial ratio $Q$ have a distorted distribution and the Lyapunov time $\tau_L$ likely has a heavy tail corresponding to highly chaotic events.[5]

Using fixed thresholds (e.g., "$\tau_L < 2$") is risky; if the data distribution is skewed, such a rule might select 50\% of the data or only 1\%. In contrast, a quantile-based rule like "select the bottom 25\% of simulations by $\tau_L$" guarantees that the resulting stage will contain exactly one-quarter of the samples it is applied to. This provides two key benefits:
\begin{enumerate}
    \item \textbf{Avoids Empty Stages:} It ensures that every defined stage will be populated with a predictable number of samples, preventing the issue of empty stages that plagued the original fixed-threshold design.[5]
    \item \textbf{Ensures Smooth Pacing:} It provides a balanced progression of difficulty. By dividing the chaotic systems into, for example, four equally sized quantile bins, the curriculum ensures the PINN is not suddenly overwhelmed by a disproportionately large number of difficult examples at a single transition. This smooth pacing is essential for stable training.[16]
\end{enumerate}
This technique is a well-established best practice in curriculum learning for regression and simulation tasks, as it automatically adapts the stage boundaries to the empirical distribution of the data, ensuring a well-proportioned and steadily advancing curriculum.[5]
